\documentclass[10pt,a4paper]{amsart}
\usepackage[margin=19mm]{geometry}
\usepackage{amsmath,amssymb,mathtools}
\usepackage[scr=rsfs,cal=euler]{mathalfa}
\usepackage{xfrac}
\usepackage[unicode,hidelinks,bookmarks=false]{hyperref}
\newcommand{\ie}{i.e.\ }
\newcommand{\cf}{cf.\ }
\newcommand{\resp}{respectively\ }
\newcommand{\Subsection}{\subsection}
\providecommand{\nobreakdash}{\nobreak\hbox{-}}
\setlength{\emergencystretch}{2em}
\multlinegap0pt
\usepackage{amssymb}
\usepackage[shortlabels]{enumitem}
\newtheorem{lemma}{Lemma}
\newcommand{\R}{\mathbb{R}}
\newcommand{\Z}{\mathbb{Z}}
\title{The Riemann $\Xi$-function from primitive Markovian cycles II: Strip rigidity and divisor identification}
\author{Douglas F. Watson}
\date{Source: arXiv:2602.06080v1 (CC BY 4.0)}
\begin{document}
\maketitle

\noindent\emph{Source and licence.} The Riemann $\Xi$-function from primitive Markovian cycles II: Strip rigidity and divisor identification. Version arXiv:2602.06080v1 (CC BY 4.0), distributed by arXiv under the Creative Commons Attribution 4.0 International licence, http://creativecommons.org/licenses/by/4.0/, as recorded in the arXiv licence metadata for this exact version and shown on the abstract page https://arxiv.org/abs/2602.06080v1. Full licence notice: \texttt{licenses/arxiv-2602.06080-CC-BY-4.0.txt}.\par\smallskip

\noindent\textit{Editorial scope.} Counted unit: the article's lemma lem:ULCLT (source0.tex lines 1076--1096). The article's complete original proof of it is delivered with it (source0.tex lines 1098--1254, one \texttt{proof} environment of 157 lines). No mathematics is written, repaired or re-derived: the statement and the proof are the article's own lines, in the article's own notation, and no retained line is substituted or re-flowed.  No further source-local context is needed: every cross-reference of every retained line resolves inside the two delivered spans.  Nothing was omitted from any retained span, so the package carries no omission record.  No retained line contains a citation, so the wrapper carries no bibliography and no bibliography entry.  No retained line contains a figure environment, an image-inclusion command or drawing code, so no image is delivered and the package declares no asset dependency.  Editorial formatting only, in addition: an AMS \texttt{amsart} preamble that restates the article's own declarations and macros used by the retained lines, namely the environments \texttt{lemma} on separate counters, together with the standard packages those lines need.  The retained environments print in this document as Lemma 1 in order of appearance, whereas the article prints the same items with the numbering of its own full document.  The complete native source of the article as distributed in the arXiv e-print for this exact version is bundled unchanged as \texttt{sources/194078/source0.tex}.
\begin{lemma}[Uniform local CLT estimate (translation-invariant cycle)]
\label{lem:ULCLT}
Assume the conductances are translation-invariant, so that the generator on $\Z/N\Z$ is
\[
(\mathcal L_N f)(j)=a\bigl(f(j+1)-f(j)\bigr)+a\bigl(f(j-1)-f(j)\bigr),
\]
and the effective diffusion constant is $D=a$. Let $p_t^{(N)}(j)$ denote the transition
kernel started from $0$ and viewed on $\Z$ via the canonical lift. Then for every compact
interval $[t_0,t_1]\subset(0,\infty)$ there exists $C=C(t_0,t_1)$ such that
\[
\sup_{t\in[t_0,t_1]}\sup_{j\in\Z}
\Bigl|
p_t^{(N)}(j)
-
\frac{1}{\sqrt{4\pi Dt}}\exp\!\Bigl(-\frac{j^2}{4Dt}\Bigr)
\Bigr|
\le
\frac{C}{N},
\]
for all sufficiently large $N$.
\end{lemma}

\begin{proof}
Because the chain is translation-invariant on $\Z/N\Z$, the Fourier modes
$e_k(j):=e^{2\pi i k j/N}$ diagonalize $\mathcal L_N$, and the corresponding eigenvalues are
\[
\lambda_k
=
a\bigl(e^{2\pi i k/N}-1\bigr)+a\bigl(e^{-2\pi i k/N}-1\bigr)
=
-2a\Bigl(1-\cos\frac{2\pi k}{N}\Bigr)
\qquad (k=0,1,\dots,N-1).
\]
Hence the heat kernel on the cycle has the exact representation
\begin{equation}\label{eq:fourier-heatkernel}
p_t^{(N)}(j)
=
\frac{1}{N}\sum_{k=0}^{N-1}\exp\!\bigl(t\lambda_k\bigr)\,e^{2\pi i k j/N}.
\end{equation}

Fix $[t_0,t_1]\subset(0,\infty)$. We compare \eqref{eq:fourier-heatkernel} to the Gaussian
kernel on $\R$ by isolating the small-frequency contribution and bounding the tail.

\smallskip
\noindent\textbf{Step 1: quadratic approximation of the dispersion relation.}
Write $\theta_k:=2\pi k/N\in[0,2\pi)$. For $|\theta|\le \pi$ one has the Taylor expansion
\[
1-\cos\theta=\frac{\theta^2}{2}+O(\theta^4),
\]
with an absolute implied constant. Therefore, for $|k|\le N/4$ (so that $|\theta_k|\le \pi/2$),
\begin{equation}\label{eq:dispersion-approx}
t\lambda_k
=
-2a t\Bigl(1-\cos\theta_k\Bigr)
=
-a t\,\theta_k^2 + O\!\Bigl(t\,\theta_k^4\Bigr)
=
-4\pi^2 D t\,\frac{k^2}{N^2}+O\!\Bigl(t\,\frac{k^4}{N^4}\Bigr),
\end{equation}
uniformly for $t\in[t_0,t_1]$.

\smallskip
\noindent\textbf{Step 2: tail bound for large frequencies.}
For $|k|\ge N/4$ (interpreting $k$ modulo $N$ and taking the representative in
$\{-\lfloor N/2\rfloor,\dots,\lfloor (N-1)/2\rfloor\}$), we have $|\theta_k|\in[\pi/2,\pi]$ and
hence $1-\cos\theta_k\ge c_0$ for an absolute constant $c_0>0$. Thus
\[
\Re(t\lambda_k)\le -2a t\,c_0 \le -2a t_0 c_0,
\]
so the tail contribution satisfies, uniformly in $t\in[t_0,t_1]$ and $j\in\Z$,
\begin{equation}\label{eq:tail}
\Bigl|\frac{1}{N}\sum_{|k|\ge N/4}\exp(t\lambda_k)\,e^{2\pi i k j/N}\Bigr|
\le
\frac{1}{N}\sum_{|k|\ge N/4}e^{-2a t_0 c_0}
\le
e^{-2a t_0 c_0}.
\end{equation}
Since the right-hand side is a fixed number in $(0,1)$, it is in particular $O(1/N)$ for
all sufficiently large $N$ (absorbing the threshold into the phrase ``sufficiently large'').

\smallskip
\noindent\textbf{Step 3: replacing the low-frequency sum by an integral.}
Restrict to $|k|<N/4$ and use \eqref{eq:dispersion-approx}. Write
\[
A_{t,N}(k):=\exp\!\Bigl(-4\pi^2 D t\,\frac{k^2}{N^2}\Bigr),
\qquad
B_{t,N}(k):=\exp(t\lambda_k).
\]
Then \eqref{eq:dispersion-approx} gives
\[
B_{t,N}(k)=A_{t,N}(k)\,\exp\!\Bigl(O\!\bigl(t\,k^4/N^4\bigr)\Bigr)
=
A_{t,N}(k)\Bigl(1+O\!\bigl(k^4/N^4\bigr)\Bigr),
\]
uniformly for $t\in[t_0,t_1]$ and $|k|<N/4$. Hence
\begin{align*}
\frac{1}{N}\sum_{|k|<N/4}\bigl|B_{t,N}(k)-A_{t,N}(k)\bigr|
&\le
\frac{C_1}{N}\sum_{|k|<N/4}A_{t,N}(k)\,\frac{k^4}{N^4}\\
&\le
\frac{C_1}{N^5}\sum_{k\in\Z}k^4\exp\!\Bigl(-4\pi^2 D t_0\,\frac{k^2}{N^2}\Bigr).
\end{align*}
Estimating the last sum by a Riemann integral with the change of variables $u=k/N$ gives
\[
\sum_{k\in\Z}k^4 e^{-c k^2/N^2}
\le
C_2 N^5
\quad\text{for }c=4\pi^2Dt_0,
\]
so the preceding display is $\le C_3/N$ uniformly in $t\in[t_0,t_1]$.

Therefore, uniformly in $t\in[t_0,t_1]$ and $j\in\Z$,
\begin{equation}\label{eq:lowfreq-replace}
\Bigl|
\frac{1}{N}\sum_{|k|<N/4}e^{t\lambda_k}e^{2\pi i k j/N}
-
\frac{1}{N}\sum_{|k|<N/4}e^{-4\pi^2 D t\,k^2/N^2}\,e^{2\pi i k j/N}
\Bigr|
\le
\frac{C_3}{N}.
\end{equation}

\smallskip
\noindent\textbf{Step 4: Poisson summation and the Gaussian kernel.}
Extend the truncated Gaussian sum to all $k\in\Z$ at cost $O(1/N)$, since
\[
\frac{1}{N}\sum_{|k|\ge N/4}e^{-4\pi^2Dt\,k^2/N^2}
\le
\frac{1}{N}\sum_{|k|\ge N/4}e^{-4\pi^2Dt_0\,k^2/N^2}
\le
\frac{C_4}{N}.
\]
Thus,
\[
\frac{1}{N}\sum_{|k|<N/4}e^{-4\pi^2 D t\,k^2/N^2}e^{2\pi i k j/N}
=
\frac{1}{N}\sum_{k\in\Z}e^{-4\pi^2 D t\,k^2/N^2}e^{2\pi i k j/N}
+O\!\Bigl(\frac{1}{N}\Bigr).
\]
By the Poisson summation formula (or the standard Jacobi theta identity),
\[
\frac{1}{N}\sum_{k\in\Z}e^{-4\pi^2 D t\,k^2/N^2}e^{2\pi i k j/N}
=
\frac{1}{\sqrt{4\pi Dt}}
\sum_{m\in\Z}\exp\!\Bigl(-\frac{(j+mN)^2}{4Dt}\Bigr).
\]
Since $t\in[t_0,t_1]$ is bounded away from $0$, the terms with $m\neq0$ are exponentially small in
$N$ uniformly in $j\in\Z$, so
\[
\sup_{t\in[t_0,t_1]}\sup_{j\in\Z}
\Bigl|
\frac{1}{\sqrt{4\pi Dt}}
\sum_{m\in\Z}e^{-(j+mN)^2/(4Dt)}
-
\frac{1}{\sqrt{4\pi Dt}}e^{-j^2/(4Dt)}
\Bigr|
\le
e^{-c_5 N^2}
\le
\frac{C_5}{N}
\]
for all sufficiently large $N$.

\smallskip
\noindent\textbf{Conclusion.}
Combining the tail estimate \eqref{eq:tail}, the low-frequency replacement
\eqref{eq:lowfreq-replace}, and the Poisson-summation identification above yields
\[
\sup_{t\in[t_0,t_1]}\sup_{j\in\Z}
\Bigl|
p_t^{(N)}(j)
-
\frac{1}{\sqrt{4\pi Dt}}\exp\!\Bigl(-\frac{j^2}{4Dt}\Bigr)
\Bigr|
\le
\frac{C}{N},
\]
for a constant $C=C(t_0,t_1)$ and all sufficiently large $N$, as claimed.
\end{proof}

\end{document}
